\documentclass[11pt,a4paper]{article}
\usepackage{fontspec}
\IfFontExistsTF{Times New Roman}{\setmainfont{Times New Roman}}{\setmainfont{Tinos}}
\IfFontExistsTF{Consolas}{\setmonofont{Consolas}[Scale=0.90]}{\setmonofont{Liberation Mono}[Scale=0.90]}
\usepackage[english]{babel}
\usepackage[top=1.35cm,bottom=1.35cm,left=1.65cm,right=1.65cm]{geometry}
\usepackage{xcolor,enumitem,listings,booktabs,array,amsmath,amssymb,titlesec,fancyhdr,parskip,tabularx,lastpage,hyperref}
\usepackage{tcolorbox}
\tcbuselibrary{skins}
\definecolor{primaryblue}{RGB}{18,65,115}
\definecolor{secondaryteal}{RGB}{14,102,110}
\definecolor{CodeBg}{RGB}{248,250,252}
\definecolor{DarkGray}{RGB}{55,55,55}
\definecolor{boxbg}{RGB}{248,250,252}
\setlength{\parindent}{0pt}\setlength{\parskip}{0pt}
\setlist[itemize]{leftmargin=1.3em,itemsep=1pt,topsep=1pt,parsep=0pt,partopsep=0pt}
\setlist[enumerate]{leftmargin=1.4em,itemsep=1pt,topsep=1pt,parsep=0pt,partopsep=0pt}
\setlist[enumerate,2]{topsep=0.5pt,itemsep=0pt,parsep=0pt,partopsep=0pt}
\lstset{language=Python,basicstyle=\ttfamily\scriptsize,keywordstyle=\color{blue!75!black}\bfseries,stringstyle=\color{red!65!black},commentstyle=\color{primaryblue!80!black}\itshape,showstringspaces=false,columns=fullflexible,keepspaces=true,breaklines=true,frame=single,rulecolor=\color{primaryblue!45},backgroundcolor=\color{CodeBg},xleftmargin=0.2em,xrightmargin=0.2em,aboveskip=0.2em,belowskip=0.2em,tabsize=4}
\pagestyle{fancy}\fancyhf{}\setlength{\headheight}{14pt}
\fancyhead[L]{\footnotesize\textbf{DSAI1003 -- Fundamental Programming Concepts in Python}}
\fancyhead[R]{\footnotesize\textbf{Sets \& Dictionaries Exam -- 60 Minutes}}
\fancyfoot[L]{\scriptsize Faculty of Data Science and Artificial Intelligence -- National Economics University}
\fancyfoot[R]{\footnotesize Page \thepage\ of \pageref{LastPage}}
\renewcommand{\headrulewidth}{0.4pt}\renewcommand{\footrulewidth}{0.3pt}
\titleformat{\section}{\large\bfseries\color{primaryblue}}{}{0pt}{}
\titleformat{\subsection}{\normalsize\bfseries\color{primaryblue}}{}{0pt}{}
\titlespacing*{\section}{0pt}{3pt}{1pt}\titlespacing*{\subsection}{0pt}{2pt}{1pt}
\newtcolorbox{headerbox}{colback=white,colframe=primaryblue,boxrule=0.7pt,arc=1.5mm,left=2.5mm,right=2.5mm,top=1mm,bottom=1mm}
\newtcolorbox{answerbox}[2][]{enhanced,colback=white,colframe=DarkGray!45,boxrule=0.5pt,arc=1mm,left=2mm,right=2mm,top=1mm,bottom=1mm,title=\footnotesize\color{DarkGray}\textbf{#2},#1}
\hypersetup{colorlinks=true,linkcolor=primaryblue,urlcolor=primaryblue,pdftitle={DSAI1003 - Sets and Dictionaries Progress Exam},pdfauthor={Faculty of Data Science and Artificial Intelligence - NEU}}
\begin{document}
% PAGE 1: institutional header, answer sheet, Q1-Q5
\noindent
\begin{minipage}[t]{0.63\textwidth}\raggedright\small
\textbf{NATIONAL ECONOMICS UNIVERSITY}\\
\textbf{COLLEGE OF TECHNOLOGY}\\
\textbf{FACULTY OF DATA SCIENCE \& ARTIFICIAL INTELLIGENCE (FDA)}
\end{minipage}\hfill
\begin{minipage}[t]{0.35\textwidth}\raggedleft\footnotesize
\textbf{PROGRESS EXAMINATION}\\
Academic Year: 2026--2027\\
Semester: Fall 2026 \quad \textbf{DSAI1003}
\end{minipage}
\vspace{0.1em}
\begin{center}
{\Large\bfseries\color{primaryblue} FUNDAMENTAL PROGRAMMING CONCEPTS IN PYTHON}\\[0.05em]
{\normalsize\bfseries PROGRESS EXAM: SETS AND DICTIONARIES (CHAPTER 8)}\\[0.05em]
{\normalsize\textbf{Paper Code:} \rule{1.5cm}{0.4pt} \quad | \quad \textbf{Duration: 60 Minutes} (Excluding distribution)}
\end{center}
\begin{headerbox}
\begin{tabularx}{\linewidth}{@{}X X@{}}
\textbf{Candidate Full Name:} \dotfill & \textbf{Student ID (MSSV):} \dotfill \\[0.15em]
\textbf{Class / Cohort:} \dotfill & \textbf{Exam Room / Seat No.:} \dotfill
\end{tabularx}
\vspace{1.5mm}
\centering\setlength{\tabcolsep}{4pt}\renewcommand{\arraystretch}{1.05}
\begin{tabular}{|l|c|c|c|c||c|c|}\hline
\textbf{Section} & \textbf{A} & \textbf{B} & \textbf{C} & \textbf{D} & \textbf{Total (100)} & \textbf{Examiner Signature} \\ \hline
\textbf{Max Score} & 20 & 25 & 25 & 30 & \textbf{100} & \\[0.2em] \cline{1-6}
\textbf{Score} & & & & & & \\[0.35em] \hline
\end{tabular}
\vspace{1mm}\raggedright\scriptsize
\textbf{Instructions:} Closed-book examination; no electronic devices or Internet access. Answer all questions. Unless otherwise specified, assume Python 3 and standard runtime behavior. Built-in functions and standard collections are permitted; third-party packages (such as NumPy, Pandas) are strictly prohibited. For program-design tasks, write modular functions with clean interfaces and a \texttt{main()} demonstration. Do not detach pages.
\end{headerbox}
\vspace{1mm}
\textbf{MULTIPLE-CHOICE ANSWER SHEET} \textit{(Write A, B, C, or D in each cell.)}
\begin{center}\setlength{\tabcolsep}{5pt}\renewcommand{\arraystretch}{1.1}
\begin{tabular}{|c|c|c|c|c|c|c|c|c|c|c|}\hline
\textbf{Question} & \textbf{1} & \textbf{2} & \textbf{3} & \textbf{4} & \textbf{5} & \textbf{6} & \textbf{7} & \textbf{8} & \textbf{9} & \textbf{10} \\ \hline
\textbf{Answer} & & & & & & & & & & \\[1.5mm] \hline
\end{tabular}\end{center}
\vspace{-2.5mm}
\section*{Section A: Set \& Dictionary Foundations \hfill \normalsize\normalfont(20 points)}
\textit{Answer all 10 questions. Each correct answer is worth 2 points. Suggested time: 12 minutes.}
\begin{enumerate}[label=\textbf{Q\arabic*.},leftmargin=2em,itemsep=1.5pt,topsep=1pt]
\item Which statement correctly instantiates an empty set in Python?
\begin{enumerate}[label=\Alph*.,leftmargin=2em]
\item \texttt{s = \{\}} \item \texttt{s = set()} \item \texttt{s = []} \item \texttt{s = set(\{\})}
\end{enumerate}
\item Which of the following elements CANNOT be added to a Python \texttt{set}?
\begin{enumerate}[label=\Alph*.,leftmargin=2em]
\item \texttt{(10, 20, 30)}
\item \texttt{"data\_science"}
\item \texttt{[1, 2, 3]}
\item \texttt{frozenset([1, 2])}
\end{enumerate}
\item What occurs when executing \texttt{s.remove(99)} versus \texttt{s.discard(99)} when \texttt{99} is not present in set \texttt{s}?
\begin{enumerate}[label=\Alph*.,leftmargin=2em]
\item Both statements raise \texttt{KeyError}.
\item Both statements execute silently without raising any exception.
\item \texttt{remove()} raises \texttt{KeyError}, while \texttt{discard()} executes silently without error.
\item \texttt{remove()} returns \texttt{False}, while \texttt{discard()} raises \texttt{ValueError}.
\end{enumerate}
\item Given sets \texttt{A = \{1, 2, 3, 4\}} and \texttt{B = \{3, 4, 5, 6\}}, what is the result of \texttt{A \textasciicircum\ B}?
\begin{enumerate}[label=\Alph*.,leftmargin=2em]
\item \texttt{\{1, 2, 3, 4, 5, 6\}}
\item \texttt{\{3, 4\}}
\item \texttt{\{1, 2\}}
\item \texttt{\{1, 2, 5, 6\}}
\end{enumerate}
\item Which expression evaluates to \texttt{True} when comparing sets \texttt{A = \{1, 2\}} and \texttt{B = \{1, 2, 3\}}?
\begin{enumerate}[label=\Alph*.,leftmargin=2em]
\item \texttt{A > B}
\item \texttt{A <= B}
\item \texttt{A.issuperset(B)}
\item \texttt{len(A) == len(B)}
\end{enumerate}
\end{enumerate}
\newpage
% PAGE 2: Q6-Q10 and Q11
\begin{enumerate}[label=\textbf{Q\arabic*.},leftmargin=2em,itemsep=1.5pt,topsep=1pt,resume]
\item Which data type is INVALID as a key in a Python dictionary?
\begin{enumerate}[label=\Alph*.,leftmargin=2em]
\item \texttt{("user", 101)}
\item \texttt{"transaction\_id"}
\item \texttt{["admin", 202]}
\item \texttt{42}
\end{enumerate}
\item What does \texttt{scores.get("Bob", 0)} return if \texttt{"Bob"} is NOT a key in dictionary \texttt{scores}?
\begin{enumerate}[label=\Alph*.,leftmargin=2em]
\item It raises a \texttt{KeyError}.
\item It automatically creates \texttt{scores["Bob"] = 0} and returns \texttt{0}.
\item It returns the default value \texttt{0} without modifying \texttt{scores}.
\item It returns \texttt{None} and prints a warning.
\end{enumerate}
\item When iterating using \texttt{for x in my\_dict:}, what does variable \texttt{x} represent at each step?
\begin{enumerate}[label=\Alph*.,leftmargin=2em]
\item The values stored in \texttt{my\_dict}
\item The keys of \texttt{my\_dict}
\item A tuple of \texttt{(key, value)}
\item An integer index \texttt{0, 1, 2, ...}
\end{enumerate}
\item Which statement deletes key \texttt{'stock'} from dictionary \texttt{d} and returns its value, or \texttt{0} if missing?
\begin{enumerate}[label=\Alph*.,leftmargin=2em]
\item \texttt{del d['stock']}
\item \texttt{d.pop('stock', 0)}
\item \texttt{d.remove('stock', 0)}
\item \texttt{d.discard('stock')}
\end{enumerate}
\item Why is membership testing \texttt{x in collection} significantly faster when \texttt{collection} is a \texttt{set} or \texttt{dict} than when it is a \texttt{list} of size $n$?
\begin{enumerate}[label=\Alph*.,leftmargin=2em]
\item Sets and dictionaries use hash tables offering average $O(1)$ lookup time, whereas lists require $O(n)$ linear scan.
\item Sets and dictionaries store elements sorted sequentially, allowing binary search in $O(\log n)$.
\item Sets and dictionaries reside in CPU registers directly.
\item Lists require sorting before any membership testing.
\end{enumerate}
\end{enumerate}
\section*{Section B: Output Tracing, Dictionary Mechanics \& Diagnostics \hfill \normalsize\normalfont(25 points)}
\textit{Answer Q11--Q15. Each question is worth 5 points. Suggested time: 15 minutes.}
\subsection*{Q11. Exact Output of a Set \& Dictionary Course Roster (5 points)}
Predict the \textbf{exact output} of the following Python program, including spacing and punctuation.
\begin{lstlisting}
roster = {}
records = [
    ("CS101", "Alice"), ("MA201", "Bob"),
    ("CS101", "Charlie"), ("CS101", "Alice"),
    ("MA201", "Alice"), ("MA201", "Bob")
]
for course, student in records:
    if course not in roster:
        roster[course] = set()
    roster[course].add(student)

for c in sorted(roster.keys()):
    print(f"{c}: {len(roster[c])} students")

both = roster["CS101"] & roster["MA201"]
print("In both courses:", sorted(both))
\end{lstlisting}
\begin{answerbox}[height=3.1cm]{Output Box for Q11}\end{answerbox}
\newpage
% PAGE 3: Q12-Q15
\subsection*{Q12. Evaluate Set and Dictionary Expressions (5 points)}
Write the exact evaluated Python result for each independent expression (1 point each):
\begin{enumerate}[label=\alph*),leftmargin=2em,itemsep=3pt]
\item After \texttt{s = \{x \% 3 for x in range(10)\}}, the value of \texttt{len(s)}: \hrulefill
\item The exact value of \texttt{\{1, 2, 3\} - \{2, 3, 4\}}: \hrulefill
\item After \lstinline|d = {"a": 10, "b": 20}|, the value of \texttt{d.get("c", 99)}: \hrulefill
\item The value of \texttt{list(dict.fromkeys(["kiwi", "apple", "kiwi", "pear"]))}: \hrulefill
\item The exception class raised by \texttt{info = \{"id": 101\}; print(info["name"])}: \hrulefill
\end{enumerate}
\subsection*{Q13. Diagnose and Correct Three Set/Dictionary Defects (5 points)}
For each defect, state the underlying cause and write a corrected Python statement.\par
\renewcommand{\arraystretch}{1.22}\small
\begin{tabularx}{\linewidth}{|p{5.1cm}|X|X|}\hline
\textbf{Erroneous Code} & \textbf{Problem / Cause} & \textbf{Correction} \\ \hline
{\ttfamily\scriptsize item\_map = \{\}}\newline{\ttfamily\scriptsize item\_map[[1, 2]] = "Origin"}\newline{\itshape Goal: store coordinate key} & & \\[7.5mm]\hline
{\ttfamily\scriptsize counts = \{\}}\newline{\ttfamily\scriptsize counts["clicks"] += 1}\newline{\itshape Goal: count page visits} & & \\[7.5mm]\hline
{\ttfamily\scriptsize tags = \{"py", "ml"\}\}}\newline{\ttfamily\scriptsize tags.remove("ai")}\newline{\itshape Goal: safely ensure tag absent} & & \\[7.5mm]\hline
\end{tabularx}\normalsize
\subsection*{Q14. Frequency Counting \& Character Distribution Analysis (5 points)}
Predict the exact console output of the code below and explain briefly why non-alphabetic characters do not appear in the resulting dictionary.
\begin{lstlisting}
freq = {}
text = "data-analytics"
for ch in text:
    if ch.isalpha():
        freq[ch] = freq.get(ch, 0) + 1

repeated = [ch for ch, c in freq.items() if c >= 2]
print("Repeated chars:", sorted(repeated))
print("Count of 'a':", freq.get("a", 0))
print("Count of 'z':", freq.get("z", 0))
\end{lstlisting}
\begin{answerbox}[height=2.5cm]{Output and Explanation for Q14}\end{answerbox}
\subsection*{Q15. Multi-Set Basket Analysis (5 points)}
Three customers at a stationery store purchase the following sets of items:
\begin{itemize}[leftmargin=2em]
\item Customer A: \texttt{\{"Notebook", "Pen", "Highlighter"\}}
\item Customer B: \texttt{\{"Pen", "Stapler", "Desk Lamp"\}}
\item Customer C: \texttt{\{"Notebook", "Pen", "Desk Lamp", "Binder"\}}
\end{itemize}
Calculate: (1) Items bought by \textbf{all three} customers; (2) Items bought by Customer A \textbf{alone} (not by B and not by C); (3) Total number of \textbf{distinct catalog items} purchased across all three customers.
\begin{answerbox}[height=2.2cm]{Answer Box for Q15}\end{answerbox}
\newpage
% PAGE 4: design
\section*{Section C: Data Structure Design \& Modular Decomposition \hfill \normalsize\normalfont(25 points)}
\textit{Answer all four tasks. Suggested time: 15 minutes.}
\begin{tcolorbox}[colback=boxbg,colframe=secondaryteal,arc=1.5mm,boxrule=0.6pt,top=1.5mm,bottom=1.5mm]
\textbf{Problem Specification: E-Commerce Customer Profiling \& Tag Recommendation Engine}\\[0.5mm]
\small
An online retail platform records purchases as a list of 4-tuples: \texttt{(customer\_id, category, amount, tag\_set)}.
\begin{itemize}
\item \textbf{Customer Profile:} Maintains for each customer their cumulative spending and the set of all unique interest tags acquired across all their purchases: \texttt{spending[cust\_id] = total\_spent}, \texttt{profiles[cust\_id] = set\_of\_tags}.
\item \textbf{Category Tag Pool:} For each product category, maintains the universal set of all tags associated with products purchased in that category: \texttt{category\_catalog[category] = set\_of\_category\_tags}.
\item \textbf{Tag Recommendations:} For a target customer and category, recommends interest tags in that category that the customer has \textbf{not yet acquired}: $\text{recommended} = \text{category\_tags} \setminus \text{customer\_tags}$.
\end{itemize}
\textbf{Sample Transaction Dataset:}\\[-0.1em]
\begin{lstlisting}
transactions = [
    ("C101", "Electronics", 120.0, {"wireless", "audio", "gadget"}),
    ("C102", "Books",        35.0, {"fiction", "paperback"}),
    ("C101", "Electronics",  80.0, {"wireless", "charger"}),
    ("C103", "Books",        50.0, {"fiction", "hardcover"}),
    ("C102", "Electronics",  45.0, {"audio", "accessory"})
]
\end{lstlisting}
\end{tcolorbox}
\vspace{0.5mm}
\textbf{Task 1: Identify Helper Functions, Parameters and Return Types (6 points)}
For each component function, state its input parameters, return type, and purpose:
\begin{itemize}[leftmargin=1.8em,itemsep=2.5pt]
\item \texttt{build\_profiles(tx\_list)} $\rightarrow$ Return Type: \rule{3cm}{0.4pt} \quad Purpose: \dotfill
\item \texttt{build\_category\_catalog(tx\_list)} $\rightarrow$ Return Type: \rule{3cm}{0.4pt} \quad Purpose: \dotfill
\item \texttt{recommend\_category\_tags(cust\_id, category, profiles, catalog)} $\rightarrow$ Return Type: \rule{2.2cm}{0.4pt} \quad Purpose: \dotfill
\end{itemize}
\textbf{Task 2: Write the Modular Processing Logic (7 points)}
Using structured pseudocode or clear Python code, implement \texttt{recommend\_category\_tags}. Handle the edge cases where either the customer ID or the category does not exist in the recorded profiles.
\begin{answerbox}[height=3.8cm]{Modular Logic Work Area for Task 2}\end{answerbox}
\textbf{Task 3: Hand-Trace the Sample Dataset (6 points)}
Compute the final state of: (1) \texttt{profiles["C101"]}; (2) \texttt{category\_catalog["Electronics"]}; (3) The recommended tags for customer \texttt{"C102"} in category \texttt{"Electronics"}.
\begin{answerbox}[height=2.8cm]{Calculation Box for Task 3}\end{answerbox}
\textbf{Task 4: Technical Analysis \& Design Decisions (6 points)}
\begin{enumerate}[label=\alph*),leftmargin=1.8em,itemsep=2pt]
\item Why is storing customer tags in a dictionary of sets (\texttt{dict[str, set[str]]}) dramatically more efficient than a flat list of tuples \texttt{[(cust\_id, tag), ...]} when testing if a customer already possesses a tag?\\[-1pt]\rule{\linewidth}{0.4pt}
\item Explain why mutable objects (such as lists or sets) cannot be used as dictionary keys or set elements in Python, referring to the hash contract.\\[-1pt]\rule{\linewidth}{0.4pt}
\end{enumerate}
\newpage
% PAGE 5: coding problem 1
\section*{Section D: Applied Programming in Python \hfill \normalsize\normalfont(30 points)}
\textit{Write complete, modular Python 3 programs. Use helper functions, a \texttt{main()} demonstration driver, clean control flow, and built-in standard library constructs only. Suggested time: 18 minutes.}
\subsection*{Problem 1: Text Concordance \& Inverted Word Index (15 points)}
Write a complete program to construct an \textbf{inverted word index} (concordance) from a multi-line document.
\textbf{Requirements:}
\begin{enumerate}[label=\arabic*.,leftmargin=2em,itemsep=1.5pt]
\item Implement \texttt{clean\_word(raw)}: strip trailing/leading punctuation, convert to lowercase, and discard any token containing non-alphanumeric characters. Return the cleaned word, or \texttt{""} if invalid.
\item Implement \texttt{build\_concordance(lines, stopwords)}:
  \begin{itemize}
  \item Take a list of strings representing consecutive lines of text (1-based line numbering).
  \item Discard words appearing in set \texttt{stopwords} or words with \texttt{len(word) < 3}.
  \item Return a dictionary mapping each valid word to the \textbf{set of line numbers} where it occurs: \texttt{dict[str, set[int]]}.
  \end{itemize}
\item Implement \texttt{query\_terms(concordance, search\_terms)}: take a list of query words and return the \textbf{set of line numbers} that contain \textbf{ALL} search terms simultaneously (using set intersection). If any term is not in the concordance, return the empty set.
\item Implement \texttt{main()} to demonstrate indexing on the sample corpus and execute a multi-term query.
\end{enumerate}
\textbf{Sample Input Corpus:}\\[-0.1em]
\begin{lstlisting}
lines = [
    "Python is an expressive and versatile language.",       # Line 1
    "Data science with Python is powerful and popular.",    # Line 2
    "Versatile libraries make Python great for data.",      # Line 3
    "Statistics and machine learning rely on Python."       # Line 4
]
stopwords = {"is", "an", "and", "with", "for", "on"}
\end{lstlisting}
\textbf{Expected Behavior:}
\begin{itemize}
\item \texttt{concordance["python"]} $\rightarrow$ \texttt{\{1, 2, 3, 4\}}
\item \texttt{concordance["data"]} $\rightarrow$ \texttt{\{2, 3\}}
\item \texttt{query\_terms(concordance, ["python", "data", "versatile"])} $\rightarrow$ \texttt{\{3\}}
\end{itemize}
\begin{answerbox}[height=8.5cm]{Python Solution for Problem 1}\end{answerbox}
\newpage
% PAGE 6: coding problem 2
\subsection*{Problem 2: Cohort Skill Matrix \& Project Team Allocation (15 points)}
Write a Python program to analyze student technical skill profiles for capstone project team matching.
\textbf{Requirements:}
\begin{enumerate}[label=\arabic*.,leftmargin=2em,itemsep=1.5pt]
\item Represent the student cohort as a dictionary mapping each student's name to a set of verified skill tags: \texttt{dict[str, set[str]]}.
\item Implement \texttt{find\_qualified(cohort, required\_skills)}: return a sorted list of names of students who possess \textbf{all} skills in the set \texttt{required\_skills} (i.e., \texttt{required\_skills.issubset(student\_skills)}).
\item Implement \texttt{skill\_gaps(cohort, student\_name, required\_skills)}: return the set of skills that the student is \textbf{missing} for the project ($\text{required} \setminus \text{student}$). Return \texttt{required\_skills} if the student is not in the cohort.
\item Implement \texttt{cohort\_skill\_summary(cohort)}: return a 2-tuple \texttt{(universal\_skills, all\_unique\_skills)}:
  \begin{itemize}
  \item \texttt{universal\_skills}: the set of skills possessed by \textbf{every single} student in the cohort (intersection of all sets). If the cohort is empty, return an empty set.
  \item \texttt{all\_unique\_skills}: the set of all unique skills present in the entire cohort (union of all sets).
  \end{itemize}
\item Implement \texttt{main()} demonstrating the functions on the sample cohort with clear, formatted output.
\end{enumerate}
\textbf{Sample Student Cohort Data:}\\[-0.1em]
\begin{lstlisting}
cohort = {
    "Minh":  {"python", "sql", "git", "statistics"},
    "Linh":  {"python", "git", "html", "css"},
    "Duc":   {"python", "sql", "git", "machine-learning"},
    "Trang": {"python", "git", "statistics", "tableau"}
}
project_reqs = {"python", "sql", "git"}
\end{lstlisting}
\textbf{Expected Behavior:}
\begin{itemize}
\item \texttt{find\_qualified(cohort, project\_reqs)} $\rightarrow$ \texttt{['Duc', 'Minh']}
\item \texttt{skill\_gaps(cohort, "Linh", project\_reqs)} $\rightarrow$ \texttt{\{'sql'\}}
\item \texttt{universal\_skills} $\rightarrow$ \texttt{\{'git', 'python'\}}
\item Total unique skills across cohort $\rightarrow$ \texttt{8}
\end{itemize}
\begin{answerbox}[height=9.2cm]{Python Solution for Problem 2}\end{answerbox}
\vfill
\begin{center}\rule{0.5\textwidth}{0.4pt}\\[1mm]
\textbf{\color{primaryblue}--- END OF EXAM PAPER ---}
\end{center}
\end{document}
